% =============================================================================
\chapter{Dyskusja}
\label{ch:discussion}
% =============================================================================

\section{Odpowiedzi na pytania badawcze}
\label{sec:answers}

\paragraph{RQ1: poprawność} Paradygmat ma co najwyżej umiarkowany wpływ na poprawność. Model osiągnął
łącznie 80\,\% poprawności za pierwszym podejściem, z rozrzutem między językami od 6/10 (Java) do
10/10 (C). Rozrzut wewnątrz paradygmatów jest równie duży jak między nimi i przy jednej próbce na
zadanie żadna różnica nie jest istotna. Najwyraźniejszy sygnał jest negatywny: paradygmat logiczny
(Prolog) ma najniższą poprawność. Nawet ten sygnał osłabiają jednak dwa defekty w zestawach testów
dla Prologu.

\paragraph{RQ2: samonaprawa} Informacja zwrotna z wykonania jest skuteczna, ale ograniczona.
Pozwoliła naprawić 9 z 14 nieudanych zadań (64\,\%), wszystkie w ciągu dwóch rund; trzecia runda
nigdy nie pomogła. Naprawa się udaje, gdy informacja zwrotna wskazuje defekt (błędy kompilatora,
niespełnione asercje z komunikatami), i utyka, gdy tego nie robi (połknięte wyjątki, defekty ukryte
w kodzie testów). Błędy formatu w odpowiedziach naprawczych zmarnowały rundy w trzech zadaniach.

\paragraph{RQ3: idiomatyczność} Tu paradygmat ma największe znaczenie. Kod funkcyjny uzyskał
najwyższe oceny (3,90), a obiektowy (3,45) i logiczny (3,40) najniższe, przy czym różnica jest
największa w samym kryterium paradygmatu (4,25 wobec 3,25). Domyślny styl modelu jest algorytmiczny
i proceduralny. Pasuje on do języków imperatywnych i dobrze przekłada się na zwinięcia
i dopasowanie wzorców, ale nie modeluje dziedziny jako obiektów ani relacji, chyba że zostanie do
tego zmuszony. Idiomatyczność języka natomiast podąża za popularnością języka, a nie za jego
paradygmatem.

\paragraph{RQ4: praktyczność} Hostowany mały model sprawia, że badanie jest tanie (około \usd{1,90}
za pełny przebieg, z czego większość przypada na sędziego) i szybkie (42 minuty czasu pracy
modelu). Główną przeszkodą praktyczną jest format odpowiedzi. Na nim poległy lokalne modele
o otwartych wagach, a nawet model hostowany dodał znaczniki bloków do 38\,\% odpowiedzi i prozę do
pięciu z nich, co spowodowało połowę wszystkich błędów kompilacji.

\section{Paradygmat czy język?}
\label{sec:paradigm-or-language}

Projekt z dwoma językami na paradygmat miał rozdzielić te dwa efekty, a wyniki pokazują, dlaczego
było to konieczne. Gdybyśmy dla paradygmatu obiektowego wybrali tylko Javę, wyglądałby on na
najsłabszy pod względem poprawności za pierwszym podejściem. Przy samym Smalltalku wyglądałby
przeciętnie. Haskell i Scala różnią się równie mocno. Przy tej skali każde twierdzenie
o poprawności dla paradygmatu jest więc w istocie twierdzeniem o wybranych językach. Inaczej jest
z idiomatycznością: Java i Smalltalk \emph{obie} uzyskują najniższe oceny za idiomy paradygmatu
(3,50 i 3,00), a Haskell i Scala \emph{oba} najwyższe (4,20 i 4,30). Tutaj efekt paradygmatu jest
spójny dla obu języków z każdej pary, co czyni go bardziej wiarygodnym ustaleniem.

Popularność języka, a więc ilość jego danych treningowych, jest czynnikiem zakłócającym, który ten
projekt ogranicza, ale nie eliminuje. Pascal, Smalltalk i Prolog to języki o małych zasobach.
Pascal i Smalltalk nie są systematycznie gorsze pod względem poprawności, ale wszystkie trzy zajmują
ostatnie miejsca w idiomatyczności języka, a Prolog i Pascal wymagały najwięcej tokenów
rozumowania. Jest to zgodne z poglądem, że model łatwiej przenosi wiedzę algorytmiczną między
językami niż konwencje konkretnego języka~\cite{cassano2024knowledge}.

\section{Dlaczego wszędzie kod proceduralny?}
\label{sec:procedural-bias}

Do proceduralnego stylu domyślnego modelu przyczyniają się prawdopodobnie trzy czynniki.
Po pierwsze, zadania są \emph{algorytmami}, a ich podręcznikowe opisy~\cite{cormen2009clrs} to
imperatywny pseudokod. Im bardziej preskryptywna specyfikacja (zadania 7--10), tym niższe oceny
idiomatyczności. Po drugie, kontrakty projektowano równolegle dla wszystkich języków, wychodząc od C.
Struktury wynikowe, wartości statusu i atomy \code{none} są wierne zadaniu, ale nie zawsze naturalne
dla paradygmatu. Po trzecie, dane treningowe języków obiektowych zawierają mnóstwo kodu w stylu klas
narzędziowych, który model odtwarza. Benchmark, który miałby mierzyć \emph{projekt} obiektowy lub
logiczny, potrzebowałby zadań, których naturalnym rozwiązaniem jest model obiektów lub relacji,
na przykład symulacji, systemów regułowych czy problemów z ograniczeniami z kilkoma oddziałującymi
na siebie bytami, a nie algorytmów z jednym punktem wejścia.

\section{Zagrożenia dla trafności}
\label{sec:threats}

\paragraph{Trafność wewnętrzna}
\begin{itemize}
  \item \emph{Jedna próbka na zadanie.} Modele Claude nie przyjmują ustawienia temperatury, więc
    każda odpowiedź jest losowaniem z rozkładu. Przy jednej próbce zmiana wyniku jednego zadania
    zmienia wynik języka o 10 punktów procentowych. System obsługuje powtórzone próbki
    (\code{SAMPLES}) i są one przystępne cenowo (\cref{sec:cost}).
  \item \emph{Defekty benchmarku.} Dwa zestawy testów dla Prologu zawierają defekty
    (\cref{sec:benchmark-defects}). Podajemy surowe liczby, więc poprawność Prologu jest zaniżona.
    Nie można wykluczyć innych, nieodkrytych defektów, zwłaszcza w ręcznie napisanych zadaniach
    1--5, których nie zweryfikowano względem rozwiązania wzorcowego w każdym języku.
  \item \emph{Nierówne specyfikacje.} Specyfikacje zadań 1--5 dostosowywano ręcznie do każdego
    języka i różnią się one długością i szczegółowością. Na przykład specyfikacja problemu
    plecakowego dla Haskella jest krótsza niż dla C (\cref{app:prompts}). Zadania 6--10
    wygenerowano ze wspólnej specyfikacji i są bardziej jednolite.
  \item \emph{Nierówna informacja zwrotna.} Programy testowe różnią się ilością wypisywanej
    diagnostyki. Program testowy w Smalltalku przechwytuje każdy wyjątek, więc samonaprawa
    w Smalltalku jest trudniejsza z przyczyn niezwiązanych z modelem.
  \item \emph{Testy przechodzące na szkielecie.} 17 testów w C przechodzi na niezaimplementowanym
    szkielecie. Nie zmieniają one żadnego werdyktu rozwiązane/nierozwiązane, ale nieznacznie
    zawyżają odsetki zaliczonych testów dla C.
\end{itemize}

\paragraph{Trafność teoretyczna}
\begin{itemize}
  \item \emph{Sędzia z tej samej rodziny.} Sędzia (\sonnet{}) i generator (\haiku{}) pochodzą
    z tej samej rodziny modeli, a sędziowie LLM mogą preferować odpowiedzi podobne do
    własnych~\cite{zheng2023judging,panickssery2024selfpreference}. Niskie oceny rozwiązań
    wzorcowych (\cref{sec:results-idiomaticity}) są zgodne z takim obciążeniem, ale go nie
    dowodzą.
  \item \emph{Brak ludzkiego punktu odniesienia.} Wykazano, że sędzia jest \emph{stabilny}, ale
    nie że jest \emph{trafny}. Próbkę rozwiązań powinni ocenić ludzcy eksperci w każdym języku.
  \item \emph{Ściśnięta skala.} Sędzia używał tylko ocen 3--5, co ogranicza rozdzielczość. Oceny
    porównawcze (szeregowanie dwóch rozwiązań) lub drobniejsza rubryka mogłyby lepiej rozróżniać
    rozwiązania.
  \item \emph{Format i poprawność są splecione.} Pięć błędów kompilacji spowodowała proza
    w odpowiedzi. Liczą się one jako błędy poprawności, choć są błędami wykonywania instrukcji.
\end{itemize}

\paragraph{Trafność zewnętrzna}
\begin{itemize}
  \item \emph{Jeden generator.} Oceniliśmy jeden mały model na jednym poziomie wysiłku. Większe
    modele, inne rodziny i modele o otwartych wagach mogą wykazywać inne wzorce; pilotaż z modelami
    lokalnymi nie dał użytecznych danych (\cref{sec:local-llms}).
  \item \emph{Dobór zadań.} Dziesięć zadań algorytmicznych to wąski wycinek programowania. Zadania
    skupione na projektowaniu systemów, wejściu/wyjściu, współbieżności czy modelowaniu dziedziny
    mogłyby faworyzować inne paradygmaty.
  \item \emph{Wersje języków.} Narzędzia są ustalone (np.\ Scala 2.11, GNU Smalltalk 3.2.5,
    \cref{tab:languages}), a idiomy nowszych wersji wykraczają poza zakres badania.
\end{itemize}

\paragraph{Trafność wnioskowania statystycznego} Podajemy statystyki opisowe i testy dokładne tylko
dla największych różnic. Przy $n = 10$ zadaniach na język i jednej próbce żadne porównanie nie
osiąga istotności na poziomie $\alpha = 0,05$. Właściwą analizą dla większego przebiegu byłyby
powtórzone próbki (pass@$k$) i model efektów mieszanych z zadaniem i językiem jako czynnikami.

\section{Wnioski dla twórców benchmarków międzyparadygmatowych}
\label{sec:lessons}

\begin{enumerate}
  \item \textbf{Unikać nasyconych problemów.} Problemy, które rozwiązuje każdy model, nie mogą
    ujawnić różnic (\cref{sec:humaneval}). Zadania z wieloma funkcjami, ścisłymi specyfikacjami
    i testami obciążeniowymi to umożliwiły.
  \item \textbf{Każdy kontrakt powinien sam się dokumentować.} Nazwy, kolejność i tryby argumentów
    muszą być widoczne w kontrakcie każdego języka, a nie tylko w nagłówku C
    (\cref{sec:benchmark-defects}).
  \item \textbf{Generować zestawy testów z jednej specyfikacji.} Ręczne przenoszenie testów między
    siedmioma językami wprowadza rozbieżności, takie jak błąd o jeden w Prologu. Umiejętność
    oparta na specyfikacji, użyta dla zadań 6--10, pozwoliła tego uniknąć.
  \item \textbf{Błędy powinny być informatywne.} Programy testowe powinny wypisywać wartości
    oczekiwane i rzeczywiste i nie mogą połykać wyjątków, ponieważ jakość informacji zwrotnej
    ogranicza samonaprawę.
  \item \textbf{Mierzyć format odpowiedzi osobno.} Należy liczyć znaczniki bloków i tekst
    niebędący kodem i oddzielać je od poprawności; inaczej porównanie mierzy wykonywanie
    instrukcji.
  \item \textbf{Walidować przed wydaniem pieniędzy.} Tryb bazowy (szkielety nie przechodzą,
    wzorce przechodzą, testy przechodzące na szkielecie są rejestrowane) nic nie kosztuje
    i wychwytuje zepsutą infrastrukturę testową przed pierwszym płatnym wywołaniem. Nie wychwyci
    jednak testu, który jest błędny w ten sam sposób dla szkieletu i dla każdego rozwiązania,
    więc rozwiązanie wzorcowe dla każdego języka pozostaje mocniejszą kontrolą.
\end{enumerate}
